\begin{vidu}
	Cho khối lượng của hạt nhân theo khối lượng nguyên tử bằng số khối của nó. Số Avogadro $N_A=6,022.10^{23}\ mol^{-1}$. Điện tích của electron là $-1,6.10^{-19}\ C$. Xét 5,6 gam \Hn{26}{56}{Fe}:
	\begin{itemize}
	\item[a)] Tính số neutron, proton, nuclôn.
	\item[b)] Tính điện tích của hạt nhân. 
	\item[c)] Có thể coi hạt nhân nguyên tử như một quả cầu bán kính $R = 1,2.10^{- 15}\sqrt[3]{A}\ m$, trong đó A là số khối. Mật độ điện tích của hạt nhân sắt  bằng bao nhiêu?
\end{itemize}
	\loigiai{
	\begin{itemize}
		\item[a)] Số nguyên tử sắt trong 5,6 g: $N_{Fe}=n.N_A=\dfrac{m}{M}.N_A=\dfrac{5,6.6,022.10^{23}}{56}=6,022.10^{22}$ nguyên tử.\\
		Vậy $\begin{cases}
		\ct{Số proton}=Z.N_{Fe}=26.6,022.10^{22}=1,566.10^{24}.\\
		\ct{Số nuclôn}=A.N_{Fe}=56.6,022.10^{22}=3,372.10^{24}.\\
		\ct{Số neutron}=(A-Z).N_{Fe}=(56-26).6,022.10^{22}=1,807.10^{24}.
		\end{cases}$ 
		\item[b)] Điện tích của hạt nhân Fe là: $Q_{Fe}=Ze=26.1,6.10^{-19}=4,16.10^{-18}\ C$.
		\item[c)] Mật độ điện tích của hạt nhân là lượng điện tích trong 1 đơn vị thể tích:
		\begin{align*}
			n = \dfrac{Z.e}{V} = \dfrac{26.1,6.10^{-19}}{\dfrac{4}{3}.\pi.\left(1,2.10^{- 15}\sqrt[3]{{56}}\right)^3} =1,026.10^{25}\ C/{m^3}.	
		\end{align*}
\end{itemize}
	}
\end{vidu}
	\begin{vidu}%[3P7-1.2K][Loai1]
	\nguon{MH - 2019} Biết số Avogadro là $6,02.10^{23}\ mol^{-1}$. Số neutron có trong 1,5 mol \Hn{3}{7}{Li} là
	\choice
	{$6,32.10^{24}$}
	{$2,71.10^{24}$}
	{$9,03.10^{24}$}
	{\True $3,61.10^{24}$}
	\loigiai{Số neutron có trong 1,5 mol \Hn{3}{7}{Li} là: 
		$N_n=\left(A-Z\right)n.N_A=\left(7-3\right).1,5.6,02.10^{23}=3,61.10^{24}$}
\end{vidu}
\begin{vidu}%[3P7-1.2K][Loai1]
	Cho khối lượng của hạt nhân theo khối lượng nguyên tử bằng số khối của nó. Số Avogadro $N_A=6,022.10^{23}\ mol^{-1}$. Số proton có trong 0,27 gam \Hn{13}{27}{Al} là
	\choice
	{\True $7,826.10^{22}$ hạt}
	{$6,826.10^{22}$ hạt}
	{$9,826.10^{22}$ hạt}
	{$8,826.10^{22}$ hạt}
	\loigiai{
		\begin{itemize}
			\item Trong hạt nhân $\Hn{13}{27}{Al}$ có $\begin{cases}
			13\ \text{proton}\\
			14\  \text{neutron}
			\end{cases}$ 
			\item Do đó, số proton là: $N_p = 13.\dfrac{0,27}{27}.{N_A} = 7,826.10^{22}$ hạt.
		\end{itemize}
	}
\end{vidu}
\begin{vidu}%[3P7-1.2B][Loai1]
	Cho khối lượng của hạt nhân theo khối lượng nguyên tử bằng số khối của nó. Số Avogadro $N_A=6,022.10^{23}\ mol^{-1}$. Trong 119 gam uranium \Hn{92}{238}{U} có số proton xấp xỉ là
	\choice
	{$4,4.10^{25}$}
	{$7,2.10^{25}$}
	{\True $2,77.10^{25}$}
	{$2,2.10^{25}$}
	\loigiai{
		Số proton có trong 119 g chất \Hn{92}{238}{U} là: $\dfrac{mZ}{A}. 6,02.10^{23}=\dfrac{119. 92}{238}. 6,02.10^{23}=2,7692.10^{25}$.
	}
\end{vidu}
\begin{vidu}%[3P7-1.2K][Loai1]
	Cho khối lượng của hạt nhân theo khối lượng nguyên tử bằng số khối của nó. Số Avogadro $N_A=6,022.10^{23}\ mol^{-1}$. Số neutron có trong 2 gam \Hn{27}{60}{Co} là 
	\choice
	{$5,254.10^{23}$ hạt}
	{$4,327.10^{23}$ hạt} 
	{$7,236.10^{23}$ hạt}
	{\True$6,622.10^{23}$ hạt}
	\loigiai{
		\begin{itemize}
			\item Một nguyên tử $\Hn{27}{60}{Co}$ có 33 neutron.
			\item 2 gam nguyên tử $\Hn{27}{60}{Co}$ có số mol là: $n = \dfrac{2}{60} = \dfrac{1}{30}$ mol.
			\item 1 mol có $N_A = 6,02.10^{23}$ nguyên tử .
			\item Do đó, trong 2 gam nguyên tử $\Hn{27}{60}{Co}$ có số neutron là $33.\dfrac{1}{30}.N_A
			= 6,622.10^{23}$.
		\end{itemize}
	}
\end{vidu}
\begin{vidu}%[3P7-1.2B][Loai1]
	Cho khối lượng của hạt nhân theo khối lượng nguyên tử bằng số khối của nó. Số Avogadro $N_A=6,022.10^{23}\ mol^{-1}$. Trong 59,50 g \Hn{92}{238}{U} có số neutron xấp xỉ là
	\choice
	{$2,38.10^{23}$}
	{\True $2,20.10^{25}$}
	{$1,19.10^{25}$}
	{$9,21.10^{24}$}
	\loigiai{
		Số neutron có trong 59,50 g chất $_{92}^{238}U$ là: $\dfrac{m\left( A-Z \right)}{A}. 6,02.10^{23}=\dfrac{59,5. 146}{238}. 6,02.10^{23}=2,20.10^{25}$.
	}
\end{vidu}
\begin{vidu}%[3P7-2.1B][Loai1]
	Biết khối lượng mol của uranium $\Hn{92}{238}{U}$ là 238 g/mol. Số Avogadro $N_A=6,022.10^{23}\ mol^{-1}$. Số neutron trong 119 gam uranium $\Hn{92}{238}{U}$ xấp xỉ là
	\choice
	{$8,8.10^{25}$ hạt}
	{$2,2.10^{25}$ hạt}
	{$1,2.10^{25}$ hạt}
	{\True $4,4.10^{25}$ hạt}
	\loigiai{
		Ta có: $\Hn{92}{238}{U} \to \begin{cases}
		A = 238\\
		Z = 92
		\end{cases} \Rightarrow N_n = A - Z = 146 \Rightarrow {N_{t\left(n \right)}} = {N_n}.\dfrac{119}{{238}}.{N_A} = 4,{4.10^{25}}$.
	}
\end{vidu}
\begin{vidu}%[3P7-1.2B][Loai1]
	Cho khối lượng của hạt nhân theo khối lượng nguyên tử bằng số khối của nó. Số Avogadro $N_A=6,022.10^{23}\ mol^{-1}$. Số nuclôn có trong 21,4 gam $\Hn{47}{107}{Ag}$ là
	\choice
	{$7,224.10^{24}$}
	{$1,6856.10^{24}$}
	{$3,3712.10^{24}$}
	{\True $1,29.10^{25}$}
	\loigiai{
		Số nuclôn có trong 21,4 g chất $_{47}^{107}Ag$ là:  $m. 6,02.10^{23}=21,4. 6,02.10^{23}=1,29.10^{25}$.
	}
\end{vidu}
\begin{vidu}%[3P7-2.1B][Loai1]
	Biết điện tích của electron là $ - 1,{6.10^{- 19}}\ C$. Điện tích của hạt nhân nguyên tử $\Hn{2}{4}{He}$ là
	\choice
	{$ - 3,2.10^{- 19}\ C$}
	{$6,4.10^{- 19}\ C$}
	{$6,4.10^{- 19}\ C$}
	{\True $3,2.10^{- 19}\ C$}
	\loigiai{
		Điện tích của hạt nhân là: $2|e| = 2.1,6.10^{- 19}\ C$.
	}
\end{vidu}
\begin{vidu*}%[3P7-1.2G][Loai1]
	Có thể coi hạt nhân nguyên tử như một quả cầu bán kính $R = 1,2.10^{- 15}\sqrt[3]{A}\ m$, trong đó A là số khối. Biết điện tích của electron là $ - 1,6.10^{- 19}\ C$. Mật độ điện tích của hạt nhân vàng $\Hn{79}{197}{Au}$ bằng bao nhiêu?
	\choice
	{\True $8,9.1{0^{24}}\ C/{m^3}$}
	{$2,3.10^{17}\ C/{m^3}$}
	{$1,8.10^{24}\ C/{m^3}$}
	{$1,2.10^{15}\ C/{m^3}$}
	\loigiai{
		Mật độ diện tích của hạt nhân là lượng điện tích trong 1 đơn vị thể tích:
		\begin{align*}
			n = \dfrac{Z.e}{V} = \dfrac{79.e}{\dfrac{4}{3}.\pi.{\left({1,2.10^{- 15}\sqrt[3]{{197}}} \right)}^3} = 8,876.19^{24}\ C/{m^3}.
	\end{align*}
}
\end{vidu*}